% Generated draft: requires mathematical and source review.
\documentclass[11pt]{article}
\usepackage[margin=25mm]{geometry}
\usepackage{fontspec}
\setmainfont{Latin Modern Roman}
\usepackage{amsmath,amssymb,mathtools,unicode-math}
\setmathfont{Latin Modern Math}
\usepackage{enumitem,needspace,xurl,hyperref}
\setlength{\parindent}{0pt}
\setlength{\parskip}{5pt}
\setlength{\emergencystretch}{3em}
\newcommand{\N}{\mathbb{N}}
\newcommand{\Z}{\mathbb{Z}}
\newcommand{\Q}{\mathbb{Q}}
\newcommand{\R}{\mathbb{R}}
\newcommand{\C}{\mathbb{C}}
\newcommand{\F}{\mathbb{F}}
\newcommand{\A}{\mathbb{A}}
\newcommand{\PP}{\mathbb{P}}
\newcommand{\E}{\mathbb{E}}
\newcommand{\eps}{\varepsilon}
\newcommand{\dd}{\,\mathrm{d}}
\begin{document}
\section*{UnsolvedMath ID: 30000040}
\textbf{Source problem code:} catalogue entry
\subsection{Definitions and mathematical statement}
The source asks how close an \(L^1\) function can be to a convolution idempotent, but does not specify the convolution algebra, ambient group, distance, admissible idempotents, or constraints on the function. Thus it does not determine one quantitative proposition.

One standard framework, stated only to identify the missing data, fixes a locally compact group \(G\) with left Haar measure and complex \(L^1(G)\), with
\[
 (f*g)(x)=\int_G f(y)g(y^{-1}x)\,\dd y,
 \qquad\mathcal I(G)=\{p\in L^1(G):p*p=p\}.
\]
For a fixed \(f\in L^1(G)\), a natural distance is
\[
 d_1(f,\mathcal I(G))=\inf_{p\in\mathcal I(G)}\|f-p\|_1.
\]
The source could ask to determine or bound this quantity, but it does not choose this interpretation. A different stability problem would compare it with the idempotence defect \(\|f*f-f\|_1\), perhaps under a bound on \(\|f\|_1\); neither that hypothesis nor a requested error function appears in the record.

The zero function is an idempotent in this framework, and functions already idempotent have distance zero. Therefore any proposed universal positive gap or nontrivial approximation restriction needs additional conditions, such as exclusion of zero, normalization, positivity, support restrictions, or a prescribed target class. These are possibilities to be specified, not assumptions imported into the source. If idempotents are instead measures, distributions, multipliers, or convolution operators, the ambient space and meaning of distance must also change. Ordinary versus twisted convolution is likewise not fixed.

The supplied literature triage dated 21 August 2026 calls the formulation underspecified, reports no located later resolution in the checked literature, and labels the record unclear/open. This is a source-limitation assessment rather than evidence for a particular unresolved stability theorem. No one of the possible completions is silently selected as the original target.
\subsection{Short English statement}
The record lacks the group, convolution, distance, and admissible-idempotent conditions needed for a unique theorem. The draft defines a standard \(L^1\)-distance framework as a possible interpretation and leaves the quantitative target explicitly unresolved.
\subsection{Sources}
\begin{itemize}
\item Local UnsolvedMath registry, permanent ID \(30000040\), exact source code \texttt{catalogue entry}, Oberwolfach Reports Open Problems. Statement status \texttt{corrected\_verified}; only \(L^1\) notation was restored and the broad wording was retained.
\item ``Mini-Workshop: Wavelets and Frames (2004),'' Oberwolfach Reports, pp. 479--544, DOI \texttt{10.4171/owr/2004/10}: \url{https://doi.org/10.4171/owr/2004/10}. Proposed year 2004.
\item Supplied literature label ``OWR 10/2004'': \url{https://doi.org/10.4171/OWR/2004/10}; supplied ``Feichtinger survey'': \url{https://doi.org/10.1007/s43670-024-00095-9}. Fuller bibliographic details and a precise target theorem are not supplied.
\item Literature checked 21 August 2026: supplied status \texttt{open}, with an explicit unclear/underspecified formulation assessment. Registry created 25 August 2026 and updated 16 September 2026. This dated triage requires independent verification before citation.
\item Registry curation attribution: Ulam AI / UnsolvedMath Contributors, CC BY 4.0. Underlying report and literature retain their supplied provenance; this statement-only draft requires human review.
\end{itemize}
\end{document}
